\chapter*{Introduction}
\addcontentsline{toc}{chapter}{Introduction} % to include the introduction to the table of content
\adjustmtc
\markboth{Introduction}{} %To redefine the section page head

Embedded systems operate under real-time constraints. They interact directly with hardware peripherals and with asynchronous external events. Both properties make them difficult to observe and to debug. Traditional techniques such as serial logging or software breakpoints are often insufficiently informative, and they are intrusive enough to perturb the behavior under examination. Modern Arm-based microcontrollers address these limitations by integrating hardware debug and trace capabilities through the Arm CoreSight architecture, which permits non-intrusive visibility into instruction execution, memory accesses and system activity.

STM32 microcontrollers expose a substantial subset of these capabilities. Despite the sophistication of the underlying hardware, the features remain locked behind proprietary tooling, or surface only as the limited workflows that commercial development environments choose to expose. Within mainstream open-source debugging environments, instruction trace is rarely treated as a first-class debugging primitive, and exists instead as a fragmented collection of low-level mechanisms requiring manual integration.

Using instruction trace in practice involves extracting raw trace data through custom scripts, decoding packets through a standalone framework such as OpenCSD, and relying on visualization or analysis tools disconnected from the debugging workflow itself. The resulting process lacks standardization, interoperability and accessibility. Advanced trace-assisted debugging is therefore placed outside the reach of many embedded developers, despite the hardware support already present on the devices they use.

This project addresses that gap through a study of the trace and debug infrastructure available on STM32 microcontrollers, spanning CoreSight components, transport protocols, debug probes, software tooling and host-side integration mechanisms. The work then proposes a practical and reproducible workflow for instruction trace capture, decoding and interpretation, designed to integrate with existing environments while remaining independent of proprietary abstractions.

The report is organized in five chapters. The first will present the host organization and will state the problem the project addresses, together with a survey of what existing tools provide. The second will establish the theory the later chapters rely upon, covering debugging fundamentals, program representation, the Arm debug architecture and the CoreSight trace architecture. The third will state the objectives, the requirements derived from them and the working environment. The fourth will describe the architecture of the solution, its decomposition into components and the reasons behind each decision. The fifth will describe the realization of the three deliverables, the tests performed and the results obtained.
